\chapter{Đào sâu: Khám phá dịch vụ --- bên trong Eureka, các kiểu lỗi và phương án Kubernetes}
\label{ch:dd-discovery}

Chương~3 giải thích Eureka trong bốn gạch đầu dòng --- đăng ký,
heartbeat, cache, tự bảo toàn --- và Chương~27 phác họa việc cho nó nghỉ
hưu trong hai dòng. Giữa hai chỗ đó là mọi thứ người vận hành thực sự
gặp: các bộ đếm thời gian quyết định một instance đã chết còn nhận lưu
lượng bao lâu, ngưỡng chuyển registry sang chế độ nó ngừng nói thật, và
sự thật khó xử rằng trên Kubernetes dự án chạy hai hệ thống khám phá
cùng lúc mà gateway chỉ hỏi một trong hai. Chương này đo những thứ đó,
làm gãy chúng, rồi chỉ ra cấu hình chính xác để gỡ Eureka mà không làm
hỏng Compose.

Ba câu hỏi sắp xếp chương này:

\begin{enumerate}
  \item \textbf{Registry giải quyết gì mà DNS không giải quyết}, và vì
  sao stack cổ điển đặt việc tra cứu ở \emph{client}?
  \item \textbf{Registry có thể cũ tới mức nào} --- tính bằng giây, với
  phép tính --- và sự cũ đó gây gì trong lúc deploy?
  \item \textbf{Cái gì thay thế từng tính năng của Eureka} khi nền tảng
  là Kubernetes, và làm sao chuyển đổi mà không có ngày ``cắt cầu dao''?
\end{enumerate}

\section{Vì sao cần registry, và vì sao ở client}

Vấn đề là địa chỉ. Instance có IP đổi theo mỗi lần khởi động lại; bên
gọi cần một cái tên ổn định. Có ba họ câu trả lời:

\begin{itemize}
  \item \textbf{Load balancer phía server.} Bên gọi quay một địa chỉ cố
  định (một nginx, một AWS ELB); nó biết các instance. Đơn giản cho bên
  gọi, nhưng mỗi bước nhảy thêm độ trễ và bộ cân bằng tự nó phải được
  cập nhật.
  \item \textbf{Registry phía client} (Eureka, Consul). Mỗi instance
  đăng ký; mỗi bên gọi tải danh sách về và cân bằng cục bộ. Không thêm
  bước nhảy, không có bộ cân bằng phải scale, bên gọi sống sót qua sự
  cố registry nhờ cache --- và mọi bên gọi phải nhúng một thư viện
  client.
  \item \textbf{DNS của nền tảng} (Service của Kubernetes). Nền tảng
  theo dõi readiness của pod và giữ danh sách endpoint của một IP ảo
  luôn cập nhật; bên gọi dùng DNS thuần và không biết gì. Không thư
  viện, không tiến trình registry, nhưng chỉ bên trong nền tảng đó.
\end{itemize}

Netflix xây cái thứ hai năm 2012 vì họ chạy trên EC2 không có DNS nền
tảng để dựa vào và hàng nghìn bên gọi JVM mang được thư viện. Spring
Cloud thừa kế nó, và dự án này thừa kế từ Spring Cloud. Lịch sử đó là
câu trả lời thành thật cho ``vì sao Eureka'': không phải vì nó hợp nhất
với một cụm Kubernetes năm 2026, mà vì stack cổ điển vẫn là thứ phần
lớn công ty Java đang chạy và phần lớn phỏng vấn Java đang hỏi.

\begin{decision}[phía client, rồi phía nền tảng]
Dự án giữ Eureka qua Compose và các milestone Kubernetes đầu tiên vì hai
lý do: Compose không có DNS gắn với readiness, nên registry là cách duy
nhất để gateway bỏ qua một service đã lên nhưng chưa sẵn sàng; và chạy
hai hệ thống song song là cách bạn học nền tảng thay thế cái gì. Khoảnh
khắc Compose không còn là đích triển khai --- nó vốn chỉ là công cụ dev
--- registry là chi phí thuần: thêm một JVM, thêm một phụ thuộc khởi
động, và một ý kiến thứ hai về pod nào còn sống, bất đồng với kubelet
tới bốn phút. Mục~\ref{sec:discovery-enhance} thực hiện cú chuyển.
\end{decision}

\section{Đọc cấu hình thật}

Server là một cổng và một annotation. Toàn bộ hành vi nằm ở các giá trị
mặc định dự án không bao giờ đặt, nên đọc phía client trước:

\begin{lstlisting}[language=yaml]
# course-service/src/main/resources/application.yml
eureka:
  client:
    service-url:
      defaultZone: http://localhost:8761/eureka/     (*@\co{1}@*)
\end{lstlisting}

\begin{lstlisting}[language=yaml]
# deploy/k8s/base/course-service.yaml (env excerpt)
- { name: EUREKA_CLIENT_SERVICEURL_DEFAULTZONE,
    value: "http://discovery-server:8761/eureka/" }  (*@\co{2}@*)
- { name: EUREKA_INSTANCE_PREFER_IP_ADDRESS, value: "true" } (*@\co{3}@*)
\end{lstlisting}

Và bên tiêu thụ, các route của gateway:

\begin{lstlisting}[language=yaml]
# configurations/api-gateway.yml (excerpt)
- id: course-service-courses
  uri: lb://course-service                          (*@\co{4}@*)
  predicates:
    - Path=/api/courses/**
  filters:
    - StripPrefix=2
\end{lstlisting}

\begin{description}
  \item[\co{1}] Địa chỉ duy nhất client được cho biết. Mọi thứ khác ---
  course-service ở đâu, auth-service ở đâu --- trở về từ URL này dưới
  dạng một tài liệu registry.
  \item[\co{2}] Trên Kubernetes \code{discovery-server} được \emph{DNS
  nền tảng} phân giải thành một ClusterIP. Chính registry được tìm thấy
  bằng cơ chế nó sắp sao chép.
  \item[\co{3}] Đăng ký IP của pod, không phải tên pod (Chương~3). Địa
  chỉ đăng ký là thứ gateway sẽ quay, bỏ qua hoàn toàn Service
  \code{course-service}.
  \item[\co{4}] \code{lb://} trao cái tên cho Spring Cloud LoadBalancer,
  nó hỏi cache cục bộ của Eureka client để lấy instance và chọn một cái
  theo vòng tròn. Service Kubernetes \code{course-service:8082} không
  bao giờ được hỏi trên đường này.
\end{description}

Điểm \co{4} là câu cần nhớ từ chương này. Trên cụm, \emph{readiness
probe bảo vệ Service của Kubernetes, và gateway không dùng Service của
Kubernetes}. Thứ duy nhất đứng giữa gateway và một pod chưa sẵn sàng là
trạng thái Eureka --- thứ pod tự báo về chính nó.

\subsection{Các bộ đếm thời gian}

Mọi thứ dưới đây là mặc định dự án thừa kế. Mỗi cái là một thuộc tính
bạn đặt được; mỗi cái có lý do cho giá trị của nó.

\begin{center}
\begin{tabular}{@{}p{0.50\linewidth}rp{0.32\linewidth}@{}}
\toprule
\textbf{Thuộc tính (\code{eureka.*})} & \textbf{Mặc định} & \textbf{Ý nghĩa} \\
\midrule
\texttt{\footnotesize instance.lease-renewal-interval-in-seconds} & 30\,s & chu kỳ heartbeat \\
\texttt{\footnotesize instance.lease-expiration-duration-in-seconds} & 90\,s & im lặng lâu thế này = hết hạn \\
\texttt{\footnotesize server.eviction-interval-timer-in-ms} & 60\,s & chu kỳ quét lease hết hạn \\
\texttt{\footnotesize server.response-cache-update-interval-ms} & 30\,s & cache chỉ đọc của server \\
\texttt{\footnotesize client.registry-fetch-interval-seconds} & 30\,s & client tải lại registry \\
\texttt{\footnotesize spring.cloud.loadbalancer.cache.ttl} & 35\,s & cache của LoadBalancer \\
\texttt{\footnotesize server.renewal-percent-threshold} & 0.85 & ngưỡng tự bảo toàn \\
\bottomrule
\end{tabular}
\end{center}

Giờ là phép tính người phỏng vấn muốn nghe.

\textbf{Instance mới nhận lưu lượng.} Spring Cloud đăng ký ngay khi web
server khởi động (client Netflix thô sẽ đợi 40\,s;
\code{EurekaAutoServiceRegistration} của Spring ép một lần đăng ký theo
yêu cầu). Cache chỉ đọc của server có thể cũ tới 30\,s, lượt tải
registry của gateway chậm tới 30\,s sau đó, và cache LoadBalancer của
nó chậm tới 35\,s sau nữa. Trường hợp xấu nhất: khoảng \textbf{95 giây}
giữa ``Tomcat đang lắng nghe'' và ``gateway biết''. Thông thường: một
nửa.

\textbf{Instance chết ngừng nhận lưu lượng.} Pod bị giết mà không tắt
êm (OOM, mất node). Lease hết hạn sau 90\,s không heartbeat; lượt quét
loại bỏ chạy mỗi 60\,s, nên tới 150\,s sau khi chết server vẫn liệt kê
nó; cộng cùng 30 + 30 + 35\,s cache: tới \textbf{245 giây} --- bốn phút
--- trong đó gateway tiếp tục quay một IP không còn trả lời. Tắt êm gửi
một DELETE lúc rời đi và bỏ qua 150\,s đầu, còn lại tới 95\,s cache cũ.

\textbf{Tự bảo toàn.} Mỗi phút server so heartbeat nhận được với
heartbeat kỳ vọng: số instance $\times$ 2 (một lần mỗi 30\,s) $\times$
0,85. Dưới vạch, nó ngừng loại bỏ bất cứ thứ gì và dashboard hiện băng
đỏ. Dự án đăng ký năm client: kỳ vọng $5 \times 2 = 10$, ngưỡng
$\lfloor 10 \times 0{,}85 \rfloor = 8$, và điều kiện là ``nhiều hơn 8''.
Giết \emph{một} service và số đếm rơi xuống 8 --- không nhiều hơn 8 ---
và registry vào chế độ tự bảo toàn. Trên stack này, một service crash
nghĩa là \emph{không} instance nào bị loại bỏ nữa cho tới khi nó quay
lại. Netflix chỉnh ngưỡng đó cho hàng trăm instance, nơi mất 15\,\%
trong một phút thật sự nghĩa là phân mảnh mạng; trên năm instance, nó
nghĩa là thứ Ba.

\begin{handson}[xem một instance chết vẫn sống]
Trong Compose, dừng course-service mà không cho nó hủy đăng ký, rồi
hỏi registry về nó một phút sau:
\begin{lstlisting}[language=bash]
$ docker compose kill course-service
$ sleep 60
$ curl -s -H 'Accept: application/json' \
    http://localhost:8761/eureka/apps/COURSE-SERVICE | \
    python -c "import sys,json; a=json.load(sys.stdin)['application']; \
      print([i['status'] for i in a['instance']])"
['UP']
\end{lstlisting}
Giờ mở \code{http://localhost:8761}. Sau lượt quét kế tiếp, instance
hoặc biến mất --- hoặc, vì bốn client còn lại gửi 8 lần gia hạn mỗi
phút, trang mang băng \textsc{emergency! eureka may be incorrectly
claiming instances are up when they're not}, và course-service giữ
\code{UP} lâu tùy bạn muốn đợi. Trong lúc đó \code{curl
localhost:8080/api/courses} trả về 503 từ gateway ở mọi lần thử rơi
vào địa chỉ đã chết.
\end{handson}

\section{Ưu và nhược, thành thật}

\begin{itemize}
  \item \textbf{Ủng hộ.} Không thêm bước nhảy mạng; bên gọi tiếp tục
  hoạt động qua sự cố registry nhờ cache; một cơ chế giống hệt trên
  laptop, trong Compose và trong cụm; một dashboard cho thấy ai đã đăng
  ký; và một instance có thể tự đánh dấu \code{OUT\_OF\_SERVICE} để rút
  lui có kiểm soát.
  \item \textbf{Phản đối, trong dự án này.} Một nguồn sự thật thứ hai về
  sự sống chậm hơn kubelet tới bốn phút; readiness probe bị bỏ qua trên
  đường của gateway (\co{4} ở trên); tự bảo toàn chỉnh cho một đội tàu
  lớn gấp trăm lần; một JVM 512\,Mi không làm gì cho lưu lượng;
  \code{prefer-ip-address} là cách chữa cháy riêng cho nền tảng mà mọi
  manifest phải mang; và mọi service biên dịch với một thư viện client
  mà upstream đang ở chế độ bảo trì.
  \item \textbf{Trung lập.} Eureka là AP theo thiết kế: nó thích trả lời
  bằng dữ liệu cũ hơn là không trả lời. Đó là lựa chọn đúng cho một
  registry --- bên gọi với danh sách hơi sai có thể thử lại; bên gọi
  không có danh sách thì không làm được gì --- và cũng là lý do mọi con
  số ở trên đều là ``tới''.
\end{itemize}

\section{Cái gì gãy}

\begin{pitfall}[rolling deploy, ba mươi giây 503]
Triệu chứng: \code{kubectl rollout restart deployment/course-service}
thành công, readiness qua, và trong nửa phút gateway trả về \code{503
Service Unavailable} trên \code{/api/courses}. Nguyên nhân: pod cũ hủy
đăng ký khi SIGTERM và thoát; pod mới đăng ký với IP mới; nhưng
LoadBalancer của gateway vẫn giữ IP cũ tới 95\,s cache. Service của
Kubernetes chuyển tức thì và không ai hỏi nó. Sửa, theo chi phí tăng
dần: một hook \code{preStop} ngủ 20--30\,s để pod cũ tiếp tục trả lời
trong lúc cache cạn (\code{terminationGracePeriodSeconds} phải lớn hơn
nó); cache ngắn hơn (\code{registry-fetch-interval-seconds: 5},
\code{spring.cloud.loadbalancer.cache.ttl: 5s}) với cái giá là tải lên
registry; hoặc một filter \code{Retry} trên gateway cho lỗi kết nối để
lần thử thứ hai rơi vào IP còn sống. Hoặc gỡ registry khỏi đường đi
(Mục~\ref{sec:discovery-enhance}) và để Service làm việc nó sinh ra để
làm.
\end{pitfall}

\begin{pitfall}[đã đăng ký, UP, và trả về lỗi]
Triệu chứng: dashboard hiện course-service \code{UP}; request thất bại
với lỗi cơ sở dữ liệu. Nguyên nhân: trạng thái Eureka do instance tự
báo, và mặc định nó là \code{UP} ngay khoảnh khắc đăng ký thành công
--- nó không hỏi actuator health. Flyway có thể đã thất bại sau khi
đăng ký, Postgres có thể đã biến mất, và instance vẫn nói \code{UP} mỗi
30\,s. Sửa: \code{eureka.client.healthcheck.enabled: true} truyền
actuator health sang trạng thái registry, để một health check
\code{DOWN} thành \code{DOWN} trong Eureka ở heartbeat kế tiếp --- 30\,s
sau, cộng cache. Readiness của Kubernetes sẽ rút pod trong 5\,s; trên
đường của gateway, nó không có cơ hội.
\end{pitfall}

\begin{pitfall}[hai replica Eureka, hai câu trả lời]
Triệu chứng: sau khi scale discovery-server lên hai replica cho tính
sẵn sàng, một số bên gọi thấy một instance mà những bên khác không
thấy. Nguyên nhân: các peer sao chép đăng ký cho nhau bất đồng bộ; một
client có \code{defaultZone} liệt kê cả hai địa chỉ nói chuyện với cái
trả lời trước và có thể đọc một replica chưa nhận được bản sao. Mỗi
server cũng đếm gia hạn độc lập, nên một cái có thể vào tự bảo toàn
trong khi cái kia loại bỏ. Sửa: liệt kê cả hai peer trong
\code{defaultZone} của mọi client và của mỗi server (chúng phải đăng ký
\emph{với nhau}), chấp nhận vài giây bất nhất là cái giá của AP, và
đừng bao giờ scale một registry bạn sắp gỡ.
\end{pitfall}

\section{Chuyện gì gãy lúc 3 giờ sáng}

Pod discovery-server chết lúc 3 giờ sáng; không gì khác chết.

\begin{itemize}
  \item \textbf{Lưu lượng gateway: không ảnh hưởng.} Eureka client của
  gateway giữ registry cuối cùng trong bộ nhớ và LoadBalancer tiếp tục
  định tuyến theo nó. Lượt tải registry thất bại mỗi 30\,s và log cảnh
  báo; cache không hết hạn. Đây là lời hứa AP được giữ.
  \item \textbf{Heartbeat: thất bại lặng lẽ.} Mọi client log một lần gia
  hạn thất bại mỗi 30\,s. Không gì khác đổi; lease là khái niệm của
  server và server đã đi.
  \item \textbf{Service khởi động lại trong lúc sự cố: vô hình.} Nó
  boot, đăng ký thất bại, thử lại ở nền --- và nó không phục vụ gì qua
  gateway cho tới khi registry trở lại và 95\,s cache cạn. Service
  Kubernetes của nó có nó làm endpoint trong vài giây; không ai dùng
  Service đó.
  \item \textbf{Service chết trong lúc sự cố: vẫn được định tuyến tới.}
  Không server, không loại bỏ. Gateway giữ IP chết cho tới khi registry
  trở lại \emph{và} đã loại bỏ nó \emph{và} cache đã quay vòng: trọn
  245\,s, tính từ lúc server trở lại.
  \item \textbf{Khi discovery-server trở lại: registry lạnh.} Nó bắt đầu
  rỗng. Trong \code{wait-time-in-ms-when-sync-empty} đầu tiên (mặc định
  5 phút) nó từ chối phục vụ một registry nó coi là chưa đầy đủ, trả về
  những gì nó có; client đăng ký lại ở heartbeat kế tiếp trong vòng
  30\,s. Hồi phục hoàn toàn: khoảng một phút sau khi pod Ready.
\end{itemize}

So với đường nền tảng: danh sách endpoint của một Service do control
plane duy trì, thứ được sao chép và mọi thành phần khác đã phụ thuộc
vào. Không có câu chuyện 3 giờ sáng riêng cho DNS Kubernetes vì nó
không có tiến trình riêng nào để mất.

\section{Nâng cấp giải pháp}
\label{sec:discovery-enhance}

Mục tiêu: trên Kubernetes, định tuyến bằng DNS của Service và để
readiness probe gác lưu lượng; trong Compose, giữ Eureka để trải nghiệm
dev không đổi. Một profile Spring tách hai bên, và config-server phục
vụ tài liệu profile chồng lên tài liệu gốc.

\subsection{Bước 1: tài liệu profile cho gateway}

\code{api-gateway.yml} gốc giữ các route \code{lb://}. Một
\code{api-gateway-k8s.yml} mới bên cạnh thay thế chúng --- cả danh
sách, vì danh sách YAML ghi đè theo chỉ số và ghi đè một phần là cái
bẫy:

\begin{lstlisting}[language=yaml]
# configurations/api-gateway-k8s.yml -- served when profile k8s is active
spring:
  cloud:
    gateway:
      routes:
        - id: auth-service
          uri: http://auth-service:8081               (*@\co{1}@*)
          predicates:
            - Path=/api/auth/**
          filters:
            - StripPrefix=2
        - id: auth-oauth2
          uri: http://auth-service:8081
          predicates:
            - Path=/oauth2/**
        - id: auth-oauth2-callback
          uri: http://auth-service:8081
          predicates:
            - Path=/login/oauth2/**
        - id: course-service
          uri: http://course-service:8082
          predicates:
            - Path=/api/courses/**,/api/students/**,/api/assignments/**,/api/enrollments/**,/api/submissions/**
          filters:
            - StripPrefix=2
        - id: dictionary-service
          uri: http://dictionary-service:8083
          predicates:
            - Path=/api/dictionary/**
          filters:
            - StripPrefix=2
      httpclient:
        connect-timeout: 2000                          (*@\co{2}@*)
        response-timeout: 10s
eureka:
  client:
    enabled: false                                     (*@\co{3}@*)
\end{lstlisting}

\begin{description}
  \item[\co{1}] Một tên Service và cổng. CoreDNS phân giải nó thành
  ClusterIP; kube-proxy cân bằng qua các endpoint Ready. Readiness probe
  giờ nằm trên đường của gateway.
  \item[\co{2}] Không có registry để hỏi, phòng thủ duy nhất của gateway
  trước một endpoint chết là connect timeout nhanh; đằng nào kube-proxy
  cũng sẽ gỡ endpoint trong vài giây.
  \item[\co{3}] Eureka client vẫn trên classpath nhưng không làm gì:
  không đăng ký, không tải, không cảnh báo heartbeat. Đây là thứ cho
  phép một image phục vụ cả hai môi trường.
\end{description}

\subsection{Bước 2: Feign không cần registry}

\code{AuthServiceClient} của course-service phân giải
\code{auth-service} qua LoadBalancer. Cho nó một URL cố định tùy chọn:

\begin{lstlisting}[language=Java]
// course-service/.../client/AuthServiceClient.java
@FeignClient(
        name = "auth-service",
        url = "${app.clients.auth-service:}",   // empty -> use lb://
        fallback = AuthServiceFallback.class)
public interface AuthServiceClient {
    @GetMapping("/users/{id}")
    UserDto getUserById(@PathVariable Long id);
}
\end{lstlisting}

Khi \code{url} phân giải thành chuỗi rỗng, Feign quay về phân giải cân
bằng tải theo \code{name}; khi nó phân giải thành
\code{http://auth-service:8081}, LoadBalancer bị bỏ qua. Tài liệu
profile của course-service đặt nó:

\begin{lstlisting}[language=yaml]
# configurations/course-service-k8s.yml
app:
  clients:
    auth-service: http://auth-service:8081
eureka:
  client:
    enabled: false
\end{lstlisting}

\subsection{Bước 3: các manifest}

Diff so với \code{deploy/k8s/base/api-gateway.yaml}; mọi manifest
service đổi cùng cách:

\begin{lstlisting}
 env:
+  - { name: SPRING_PROFILES_ACTIVE, value: "k8s" }
   - { name: SPRING_CONFIG_IMPORT,
       value: "configserver:http://config-server:8888" }
-  - { name: EUREKA_CLIENT_SERVICEURL_DEFAULTZONE,
-      value: "http://discovery-server:8761/eureka/" }
-  - { name: EUREKA_INSTANCE_PREFER_IP_ADDRESS, value: "true" }
\end{lstlisting}

Và trong \code{kustomization.yaml}:

\begin{lstlisting}
 resources:
   - config-server.yaml
-  - discovery-server.yaml
   - api-gateway.yaml
\end{lstlisting}

Compose không đụng tới: không profile, tài liệu gốc, \code{lb://},
Eureka. Cùng các image chạy ở cả hai. Khi Compose cuối cùng nghỉ hưu,
xóa \code{spring-cloud-starter-netflix-eureka-client} khỏi mọi
\code{pom.xml}, module \code{discovery-server}, và các route
\code{lb://} trong tài liệu gốc --- theo thứ tự đó, mỗi việc một commit
riêng để pipeline kiểm chứng được.

\begin{handson}[chứng minh đường nền tảng]
Sau khi áp dụng profile, xác nhận gateway không còn quay thẳng IP pod
--- nó quay Service, mà endpoint theo readiness:
\begin{lstlisting}[language=bash]
$ kubectl -n ts get endpoints course-service
NAME             ENDPOINTS         AGE
course-service   10.42.0.57:8082   3d
$ kubectl -n ts rollout restart deployment/course-service
$ while true; do curl -s -o /dev/null -w '%{http_code}\n' \
    -H "Authorization: Bearer $TOKEN" \
    http://ts.local/api/courses; sleep 1; done
200
200
200   # no 503 window: the endpoint switched with readiness
\end{lstlisting}
Chạy cùng vòng lặp trước khi đổi và đếm số 503.
\end{handson}

\section{Ngoài dự án}

Consul là registry bạn gặp ngoài JVM: một agent trên mỗi host tự chạy
health check thay vì tin heartbeat tự báo, và phục vụ kết quả qua DNS
lẫn HTTP, nên bên gọi không cần thư viện. Service mesh (Istio, Linkerd)
đi thêm một bước và đặt một proxy sidecar cạnh mỗi pod làm khám phá,
retry, timeout và mTLS, để ứng dụng chỉ còn
\code{http://course-service}. Trong Spring,
\code{spring-cloud-kubernetes} cho phép tên \code{lb://} phân giải qua
API Kubernetes thay vì Eureka --- hữu ích trong lúc chuyển đổi, không
cần thiết khi chuyển xong. Và AWS Cloud Map, hậu duệ được quản lý của ý
tưởng Eureka, tồn tại chủ yếu cho ECS, nơi không có kube-proxy để khiến
nó thừa.

\section{Câu hỏi từ phỏng vấn thật}

\subsection*{Q: Gateway của bạn vẫn gửi lưu lượng tới một instance đã
chết chín mươi giây trước. Vì sao, và cắt ngắn bằng cách nào?}

Vì không gì trên đường đi nhận ra. Eureka biết một cái chết không êm
chỉ qua heartbeat thiếu: 90\,s hết hạn lease, rồi một lượt quét chạy
mỗi 60\,s, rồi 30\,s cache server, 30\,s tải của client, và 35\,s cache
LoadBalancer --- tới bốn phút đầu-cuối, và nếu registry đã trượt vào
tự bảo toàn, mãi mãi. Cắt phần đuôi trước: rút
\code{registry-fetch-interval-seconds} và TTL LoadBalancer xuống 5\,s,
tốn lưu lượng registry mà bạn kham được dưới vài trăm instance. Rồi
cắt phần đầu: làm cái chết trở nên êm (xử lý SIGTERM,
\code{terminationGracePeriodSeconds} đủ dài) để DELETE được gửi đi và
150\,s logic lease không bao giờ chạy. Rồi làm bên gọi bền bất kể: một
filter \code{Retry} cho lỗi kết nối, để một địa chỉ cũ tốn một lần thử
lại, không phải một lỗi. Câu trả lời senior kết bằng ``hoặc ngừng định
tuyến bằng registry trên Kubernetes, vì kube-proxy gỡ endpoint trong
vài giây sau khi readiness probe thất bại''.

\subsection*{Q: Làm sao một instance mới bắt đầu nhận lưu lượng trong
rolling deploy với không lỗi nào?}

Hai điều kiện và một chồng lấn. Instance mới không được nhận lưu lượng
trước khi phục vụ được: trên Kubernetes đó là readiness probe, trên
Eureka đó là chỉ đăng ký sau khi context lên hẳn --- Spring làm thế,
nhưng một health check qua trước khi Flyway xong thì không, nên hãy
probe thứ gì phụ thuộc cơ sở dữ liệu. Instance cũ phải tiếp tục phục vụ
tới khi mọi bên gọi ngừng gửi: trên Kubernetes endpoint bị gỡ lúc
SIGTERM nhưng request đang chạy hoàn tất trong grace period, và một
\code{preStop} ngủ che cho bên gọi có cache cũ. Chồng lấn là
\code{maxSurge: 1, maxUnavailable: 0}: đưa cái mới lên, đợi Ready, rồi
hạ cái cũ. Lỗi trong rollout hầu như luôn là một trong ba thứ --- probe
nói dối, cache sống lâu hơn pod, hay grace period ngắn hơn request chậm
nhất --- và bạn gọi tên được cái nào từ hình dạng của lỗi.

\subsection*{Q: Eureka so với DNS Kubernetes. Một phút.}

Cả hai ánh xạ một tên tới một tập địa chỉ; chúng khác ở ai duy trì tập
đó và ai tra cứu. Eureka: instance tự đăng ký, tự báo sức khỏe, và mọi
bên gọi mang một client tải và cache danh sách --- không cần nền tảng,
chạy trên VM, AP với vài phút cũ, và mọi bên gọi là một JVM có thư
viện. Kubernetes: control plane duy trì endpoint của Service từ label
selector và readiness probe do chính nó chạy, kube-proxy cân bằng, bên
gọi dùng DNS và không biết gì --- vài giây cũ, không thư viện, không
tiến trình registry, nhưng chỉ trong cụm. Trên Kubernetes, Eureka là ý
kiến thứ hai, chậm hơn, về sự sống. Chỉ giữ nó nếu bạn buộc phải chạy
cùng code ở nơi nền tảng không có mặt.

\subsection*{Q: Bạn chạy hai replica Eureka và chúng bất đồng về việc
một instance còn sống không. Giờ sao?}

Hãy lường trước; thiết kế cho nó. Sao chép giữa peer là bất đồng bộ và
mỗi replica loại bỏ và tự bảo toàn theo bộ đếm riêng, nên một cửa sổ
bất đồng là cố hữu --- Eureka là AP và bất đồng là cái giá của việc trả
lời trong lúc phân mảnh. Điều quan trọng là client chịu được: mọi
client liệt kê cả hai replica, thử cái thứ hai khi thất bại, và coi
registry là gợi ý mà một connect timeout có thể phủ quyết. Nếu bất
đồng kéo dài quá vài phút, một replica đang tự bảo toàn còn cái kia
không; so ngưỡng gia hạn với số instance thực, và trên đội tàu nhỏ hạ
\code{renewal-percent-threshold} hoặc tắt hẳn tự bảo toàn. Câu trả lời
sai là cố làm registry nhất quán; câu đúng là làm bên gọi thờ ơ với
replica nào chúng đọc.

\subsection*{Q: Hai trăm service, mỗi cái tải registry mỗi ba mươi
giây. Registry đang tan chảy. Bạn đổi gì?}

Bốn trăm lượt tải mỗi phút của một tài liệu liệt kê hai trăm ứng dụng
là nút thắt đầu tiên, và câu trả lời của Eureka là tải delta: client
hỏi \code{/apps/delta} sau lần tải đầy đủ đầu tiên, mặc định bật và chỉ
hỏng khi hàng đợi delta của server bị tắt. Thứ hai, cache phản hồi chỉ
đọc tồn tại chính xác để server tuần tự hóa registry một lần mỗi
30\,s, không phải mỗi request --- giữ nó. Thứ ba, chia theo zone: cấu
hình \code{availability-zone} của Eureka cho client ưu tiên registry
cục bộ và instance cục bộ, để mỗi server gánh một phần lượt tải. Quá
một nghìn instance, câu trả lời của chính Netflix là thêm server Eureka
sau DNS, và câu trả lời ngày nay là ở cỡ đó bạn đang ở trên một nền
tảng làm khám phá hộ bạn.

\begin{quiz}
\begin{enumerate}
  \item Tính độ trễ xấu nhất trước khi gateway định tuyến tới một pod
  course-service vừa khởi động, gọi tên từng cache trong chuỗi và mặc
  định của nó.
  \item Một instance bị OOM killer giết. Liệt kê các bộ đếm phải trôi
  qua trước khi gateway ngừng quay nó, và nêu cái nào một lượt tắt êm
  sẽ bỏ qua.
  \item Với năm client đăng ký, chỉ ra vì sao mất một cái đưa registry
  vào tự bảo toàn, và nêu cấu hình một dòng thay đổi điều đó trên đội
  tàu nhỏ.
  \item Trên Kubernetes, cơ chế nào gác lưu lượng tới một pod trên
  đường của gateway hôm nay, và vì sao readiness probe không? Thay đổi
  cấu hình đơn nào đưa readiness trở lại đường đi?
  \item Sau \code{eureka.client.enabled: false} dưới profile \code{k8s},
  giải thích vì sao cùng image course-service vẫn đăng ký với Eureka
  trong Compose, và \code{url = "\${app.clients.auth-service:}"} làm gì
  ở mỗi môi trường.
  \item Phỏng vấn: ``Vì sao registry tiếp tục phục vụ một instance đã
  chết trong khi Service Kubernetes thì không?'' Trả lời theo ai quan
  sát sự sống và bằng cách nào.
\end{enumerate}
\end{quiz}
